\subsection{What Is Unique About Agent Context}
\label{subsec:what_is_unique}

Formally, agent context is a time-indexed state $C_t$ continuously updated by interaction tuples $x_t = (a_t, t_t, o_t)$ of actions, thoughts, and observations~\citep{yao2023reactsynergizingreasoningacting}. Unlike standard long-text processing, this dynamic trajectory exhibits four fundamental properties that distinguish it from static prompts:

\noindent\textbf{Property 1: Dynamic Growth.} Context is not provided all at once but expands incrementally by absorbing new interactions $x_t$ with each agent step. Unlike static documents with fully exposed structures, agentic compression operates blindly on this continuous stream, dynamically distilling history without foreseeing future states.

\noindent\textbf{Property 2: Heterogeneous Composition.} Context comprises highly distinct structured content, including natural language plans, code snippets, JSON tool outputs, screenshots, and error logs. While natural language thoughts ($t_t$) may endure lossy abstraction, actions and observations often demand strict syntactic fidelity, mandating type-aware rather than uniform compression.

\noindent\textbf{Property 3: Multi-step Dependency.} Information compressed or discarded at a certain step may only be needed several steps later. Driven by continuous streaming, the utility of any specific tuple component $e$ manifests over future steps: $\mathcal{U}(e) = \mathbb{E}\left[ \sum_{\tau = t}^{T} \gamma^{\tau - t} \cdot u_\tau(e) \right]$. Consequently, it is inherently difficult to predict future demands during compression, as the true importance of $e$ remains unobservable at time $t$.

\noindent\textbf{Property 4: Error Propagation.} Errors or omissions introduced by compression are not discovered locally. Instead, a state distortion $\delta_t = d(C_t, \tilde{C}_t)$ propagates backward along the trajectory: $C_{t+k} = U^{(k)}(\tilde{C}_t, x_{t+1:t+k})$. These errors accumulate irreversibly and are treated as factual grounding by subsequent reasoning steps, permanently shifting the action space and derailing the agent from the optimal path.

\vspace{1ex}
\noindent\textit{\textbf{$\blacktriangleright$ Why Traditional Compression Fails $\blacktriangleleft$}} Traditional methods satisfy one-off budget constraints by optimizing for immediate relevance~\citep{jiang-etal-2023-llmlingua, jiang2024longllmlinguaacceleratingenhancingllms}. For example, KV cache compression typically evicts tokens with low current attention scores~\citep{xiao2024efficientstreaminglanguagemodels, NEURIPS2024_28ab4182}. However, under the four properties above, applying such static interventions to an evolving agentic loop triggers compounding distortions, as continuous interaction defies any static definition of information importance.


\subsection{A Unified Operational Pipeline for Agent Context Compression}
\label{subsec:unified_pipeline}

% Existing methods for agent context compression are fragmented, making it difficult to compare systems or reason about where compression intervenes in execution.
We formulate a unified four-stage pipeline that decomposes compression into an operator system over the context lifecycle:

{\setlength{\abovedisplayskip}{4pt}
\setlength{\belowdisplayskip}{4pt}
\setlength{\abovedisplayshortskip}{2pt}
\setlength{\belowdisplayshortskip}{2pt}
\vspace{-.3cm}
\begin{equation}
    C_{t+1} = \underbrace{R(q_{t+1}, \mathcal{M}_{t})}_{\text{Recover}} \oplus \underbrace{M\Big(\Phi\big(S(C_{t})\big)\Big)}_{\text{Store} \leftarrow \text{Compress} \leftarrow \text{Select}} \oplus \ x_{t+1} \label{eq:pipeline}
\end{equation}
}

% This formulation operationalizes the following four chronological stages:
\begin{itemize}[leftmargin=*, parsep=0pt, topsep=1pt, itemsep=1pt]
    \item \textbf{Select ($S$):} A critical triage phase where the system determines which subsets of the context require compression and which demand exact preservation.
    \item \textbf{Compress ($\Phi$):} Specific reductive or structural transformations, ranging from simple truncation to latent encoding, are executed on the selected subset.
    \item \textbf{Store ($M$):} The resulting compressed representations are committed either to the active working memory or offloaded to an external storage state \(\mathcal{M}_t\).
    \item \textbf{Recover ($R$):} Driven by the necessity of deferred utility, relevant compressed information is selectively retrieved and decoded from \(\mathcal{M}_t\) using the current query \(q_{t+1}\).
\end{itemize}

% This four-stage formulation decouples context management from the agent's active reasoning loop. Compression thus acts as a backend state-management subsystem concerned with information density and recoverability; the next two sections make this view explicit through a formal lens and a taxonomy.

% \subsection{Why Long-Context Windows Are Not a Panacea}
% 
% A natural objection is that long-context models should render compression unnecessary. In practice, larger windows help, but they do not remove the underlying problem.
% 
% First, longer contexts do not eliminate economic constraints. Even when a model can technically accept more tokens, inference cost, latency, and memory usage still grow with context length. In interactive agent systems, these costs remain material~\citep{hooper-etal-2025-squeezed, liu2024retrievalattentionacceleratinglongcontextllm}.
% 
% Second, more context is not equivalent to better context use. Empirical studies by \citet{liu2023lostmiddlelanguagemodels} on long-context models show that information in long prompts may be underutilized, and that simply increasing the window does not guarantee improved retrieval or reasoning. The issue is therefore not only capacity, but utilization~\citep{kim-etal-2025-really, tian-etal-2025-distance}.
% 
% Third, the central concern is information quality, not token count alone. A long but noisy trajectory can still degrade planning, dilute salient constraints, and increase brittleness. Compression is thus a mechanism for improving information density, not merely for fitting within a budget~\citep{du-etal-2025-context}.
% 
% In short, a larger context window raises the ceiling; context compression raises the floor.
